%% ========================================================
\section{Термохімія та нульова енергія}
%% ========================================================


Розглянемо, як із результатів квантово-хімічного
розрахунку (зокрема з гесіана) отримують термохімічні параметри:
нульову коливальну енергію, термічні поправки, ентропію та вільну енергію.
Ці параметри \emph{відносяться до газової фази} та
розраховуються в рамках \emph{моделі ідеального газу}.

У квантово-хімічних програмах, таких як \texttt{PySCF},
припускається, що молекула є ізольованою частинкою у газовій фазі,
а її внутрішня енергія складається з кількох незалежних внесків:
\[
E_{\text{total}} = E_{\text{elec}} + E_{\text{vib}} + E_{\text{rot}} + E_{\text{trans}}.
\]
Тут $E_{\text{elec}}$ --- \emph{електронна енергія системи},
отримана з розв’язку електронного рівняння Шредінґера
при фіксованих положеннях ядер (наближення Борна–Оппенгеймера).
Залежно від рівня теорії, цей внесок може бути визначений
у межах методу Хартрі–Фока (HF), теорії функціонала густини (DFT)
або більш точних пост-ХФ підходів

Для оцінки цих ядерних внесків використовується
\emph{наближення жорсткого ротатора} (для обертання молекули)
та \emph{гармонічного осцилятора} (для коливань атомів).
Відповідно, термохімічні функції молекул —
енергія, ентальпія, ентропія та вільна енергія Гіббса —
залежать не лише від електронної будови, але й від термічних ефектів,
пов’язаних із поступальним, обертальним і коливальним рухом ядер.



%% --------------------------------------------------------
\subsection{Нульова коливальна енергія (ZPE)}
%% --------------------------------------------------------

Навіть при абсолютному нулі температури ($T=0$~K) молекула не є нерухомою.
Згідно з принципом невизначеності Гайзенберґа, ядра не можуть одночасно
мати точно визначені координати й імпульси, тому вони завжди здійснюють
квантові коливання навколо рівноважного положення.

Для гармонічного осцилятора енергії рівнів задаються:
\[
E_v = \left(n + \frac{1}{2}\right)\hbar\omega, \qquad n = 0, 1, 2, \ldots
\]
Тобто навіть у найнижчому стані ($n=0$) коливальна енергія дорівнює
$\tfrac{1}{2}\hbar\omega$.
Для багатомодової молекули (з $3N-6$ нормальних мод)
нульова коливальна енергія (Zero Point Energy, ZPE) становить:
\[
E_{\text{ZPE}} = \sum_{i=1}^{3N-6} \frac{1}{2}\hbar\omega_i.
\]

Знання ZPE необхідне для точних термохімічних розрахунків (зокрема енергій реакцій), також вона изначає \emph{ізотопні ефекти}, впливає на швидкість \emph{тунелювання} через потенційні бар’єри. Для невеликих органічних молекул становить зазвичай $5$–$15$~ккал/моль .


%% --------------------------------------------------------
\subsection{Термохімічні поправки}
%% --------------------------------------------------------

При підвищенні температури ядра можуть займати збуджені коливальні,
обертальні й поступальні стани. Це змінює середню енергію системи,
її ентальпію, ентропію та вільну енергію Гіббса.

Для кожного ступеня вільності маємо:

\begin{equation*}
    E_{\text{vib}}(T) = \sum_i \frac{\hbar\omega_i}{\exp(\hbar\omega_i/k_BT) - 1},
\end{equation*}
\begin{equation*}
    E_{\text{rot}}(T) =
    \begin{cases}
        k_BT, & \text{для лінійної молекули;} \\
        \frac{3}{2}k_BT, & \text{для нелінійної.}
    \end{cases}
\end{equation*}
\begin{equation*}
    E_{\text{trans}}(T) = \frac{3}{2}k_BT.
\end{equation*}

Тоді повна енергія системи:
\begin{equation*}
    E_{\text{total}}(T) = E_{\text{elec}} + E_{\text{ZPE}} + E_{\text{vib}}(T)
    + E_{\text{rot}}(T) + E_{\text{trans}}(T).
\end{equation*}

Ентальпія визначається як:
\begin{equation*}
    H(T) = E_{\text{total}}(T) + RT.
\end{equation*}

%% --------------------------------------------------------
\subsection{Ентропія молекули}
%% --------------------------------------------------------

Ентропія --- це міра невпорядкованості системи або кількості доступних мікростанів.
Для ідеального газу вона має такі складові:
\[
    S(T) = S_{\text{trans}} + S_{\text{rot}} + S_{\text{vib}} + S_{\text{elec}}.
\]

\paragraph{1. Поступальна ентропія:}
\[
    S_{\text{trans}} = R\left[
        \ln\!\left(\frac{(2\pi m k_BT)^{3/2}k_BT}{h^3 P}\right) + \frac{5}{2}
    \right],
\]
де $m$ --- маса молекули, $P$ --- тиск.

\paragraph{2. Обертальна ентропія:}
\[
    S_{\text{rot}} = R\left[
        \ln\!\left(
        \frac{\sqrt{\pi}T^{3/2}}{\sigma}
        \left(\frac{8\pi^2k_B}{h^2}\right)^{3/2}
        \sqrt{I_A I_B I_C}
        \right) + \frac{3}{2}
    \right],
\]
де $\sigma$ --- число симетрійних перестановок, $I_A,I_B,I_C$ --- головні моменти інерції.

\paragraph{3. Коливальна ентропія:}
\[
    S_{\text{vib}} = R \sum_i
        \left[
            \frac{\theta_i/T}{\exp(\theta_i/T) - 1}
            - \ln\!\left(1 - e^{-\theta_i/T}\right)
        \right],
\]
де $\theta_i = h\nu_i/k_B$ --- характеристична температура моди.

\paragraph{4. Електронна ентропія:}
Зазвичай нехтують, бо основний стан є єдиним при кімнатних температурах.


%% --------------------------------------------------------
\subsubsection{Вільна енергія Гіббса}
%% --------------------------------------------------------

\[
    G(T) = H(T) - TS(T)
\]

Вона визначає термодинамічну стабільність сполук
і напрямок хімічних реакцій.

%% --------------------------------------------------------
\subsection{Як це реалізовано в \texttt{PySCF}}
%% --------------------------------------------------------
У пакеті \texttt{PySCF} повний термохімічний аналіз молекули базується
на енергії електронів $E_{\text{elec}}$, до якої послідовно додаються
внески нульових коливань, теплові поправки та ентропійні терміни. Термохімічні величини обчислюються
на основі гармонічних частот, отриманих із гессіана
(матриці других похідних енергії за координатами атомів),
і додаються до електронної енергії $E_{\text{elec}}$,
отриманої з розв’язку рівнянь Хартрі–Фока або DFT.

Далі програма:
\begin{enumerate}
    \item Визначає нормальні частоти $\nu_i$ після масового масштабування гессіана.
    \item Обчислює нульову енергію коливань $E_{\text{ZPE}} = \tfrac{1}{2}\sum h\nu_i$.
    \item Розраховує термічні внески $E_{\text{vib}}(T), E_{\text{rot}}(T), E_{\text{trans}}(T)$
          за статистично-механічними формулами (як наведено вище).
    \item Обчислює ентальпію $H(T)$, ентропію $S(T)$ та вільну енергію $G(T)$.
\end{enumerate}

Всі ці розрахунки виконуються у межах
\emph{наближення гармонічного осцилятора та жорсткого ротатора},
що забезпечує добру точність для малих і середніх молекул.


%% --------------------------------------------------------
\subsection{Термохімічний приклад: \ce{H2O}}
%% --------------------------------------------------------

\inputcode[mod=inline]{h2o_thermochemistry.py}


У цій програмі основна робота термохімічного аналізу була зведена до двох ключових функцій з модуля \href{https://pyscf.org/_modules/pyscf/hessian/thermo.html}{\texttt{pysch.hessian.thermo}}:

\begin{enumerate}
  \item \inlinecode{harmonic\_analysis(mol, hess)} --- виконує гармонічний аналіз молекули на основі гесіану.
        Вона обчислює частоти нормальних коливань, нормалізовані моди, редуковані маси та нульову коливальну енергію (ZPE).
        Ця функція також враховує трансляційні та обертальні ступені свободи та дозволяє виключати їх при необхідності.
  \item \inlinecode{thermo(model, freq, temperature, pressure)} --- використовує результати гармонічного аналізу для обчислення термодинамічних величин:
        ентропії, теплоємності ($C_V$, $C_P$), внутрішньої енергії, ентальпії та вільної енергії Гіббса.
        Вона враховує внески електронної, поступальної, обертальної та коливальної частин.
\end{enumerate}

Таким чином, за рахунок цих двох функцій програма зводить складні квантово-хімічні розрахунки та гармонічну апроксимацію в зручний формат для термохімічного аналізу.
Вся <<магія>> полягає в тому, що \inlinecode{harmonic\_analysis} формує основні фізичні параметри, а \inlinecode{thermo} перетворює їх у термодинамічні величини, готові для подальшого виводу у таблиці.



%\noindent
%\textit{Коментар:}
%Більшість квантово-хімічних програм (Gaussian, ORCA, PySCF тощо)
%в автоматичному режимі обчислюють ці термохімічні функції на основі
%власних частот, отриманих із гесіану при рівноважній геометрії.
%Тому результати часто подаються у вигляді:
%\[
%E_{\text{elec}}, \quad E_{\text{elec}} + \text{ZPE}, \quad H(298), \quad G(298).
%\]
%Саме ці значення використовують у подальших термодинамічних розрахунках,
%наприклад, при оцінюванні енергії реакції чи рівноважних констант.



%%% --------------------------------------------------------
%\subsubsection{Термохімічний розрахунок для \ce{CH4}}
%%% --------------------------------------------------------
%
%
%\inputcode{ch4_thermochemistry.py}
%
%\textbf{Результати для \ce{CH4} при $298.15$~K:}
%
%\begin{center}
%    \begin{tblr}{
%        colspec = {X[l,3] X[r,2] X[l,2]},
%        row{1} = {font=\bfseries},
%        hline{1,2,Z} = {solid},
%            }
%        Величина              & Значення & Одиниці      \\
%        ZPE                   & 28.03    & ккал/моль    \\
%        $E_{\text{vib}}(T)$   & 0.10     & ккал/моль    \\
%        $E_{\text{rot}}(T)$   & 0.89     & ккал/моль    \\
%        $E_{\text{trans}}(T)$ & 0.89     & ккал/моль    \\
%        $H(T) - H(0)$         & 2.48     & ккал/моль    \\
%        $S(T)$                & 44.5     & кал/(моль·K) \\
%        $C_V$                 & 6.0      & кал/(моль·K) \\
%    \end{tblr}
%\end{center}
%
%\textit{Розрахунок: B3LYP/6-311+G(2d,p)}
%
%\textbf{Порівняння з експериментом:}
%\begin{itemize}
%    \item ZPE (експ.): 27.8 ккал/моль --- відмінна згода
%    \item $S_{298}$ (експ.): 44.5 кал/(моль·K) --- ідеальна згода
%    \item Термохімічні поправки надійні для DFT
%\end{itemize}

%%% --------------------------------------------------------
%\subsection{Ізотопні ефекти}
%%% --------------------------------------------------------
%
%Заміна ізотопу змінює частоти через зміну маси:
%\[
%    \frac{\omega_{\text{D}}}{\omega_{\text{H}}} \approx \sqrt{\frac{m_{\text{H}}}{m_{\text{D}}}} = \sqrt{\frac{1}{2}} \approx 0.707
%\]

%\subsubsection{Ізотопний ефект у \ce{H2O}/\ce{D2O}}
%
%%\inputcode{h2o_d2o_isotope_effect.py}
%
%\textbf{Порівняння \ce{H2O} та \ce{D2O}:}
%
%\begin{center}
%    \begin{tblr}{
%        colspec = {X[l,3] X[r,2] X[r,2] X[r,2]},
%        row{1} = {font=\bfseries},
%        hline{1,2,Z} = {solid},
%            }
%        Мода                & \ce{H2O}    & \ce{D2O}    & Співвідношення \\
%                            & (см$^{-1}$) & (см$^{-1}$) &                \\
%        Симетричне валентне & 3657        & 2671        & 0.730          \\
%        Деформаційне        & 1595        & 1178        & 0.738          \\
%        Антисим. валентне   & 3756        & 2788        & 0.742          \\
%    \end{tblr}
%\end{center}
%
%\textit{Експериментальні дані}
%
%\textbf{ZPE ефект:}
%\begin{itemize}
%    \item ZPE(\ce{H2O}): 13.26 ккал/моль
%    \item ZPE(\ce{D2O}): 9.66 ккал/моль
%    \item $\Delta$ZPE: 3.60 ккал/моль
%\end{itemize}
%
%\textbf{Наслідки:}
%\begin{itemize}
%    \item Дейтеровані сполуки стабільніші (нижча ZPE)
%    \item Кінетичний ізотопний ефект: $k_{\text{H}}/k_{\text{D}} \approx 7$ для розриву \ce{C-H}
%    \item Використання в механістичних дослідженнях
%\end{itemize}
